% GENERATED FILE. Edit canonical lessons, then run npm run generate:books.
% canonical-source-hash: fnv1a64:e807d95a88fc77ab
% canonical-lessons: ES-C543-reanudar, ES-C543-retrasar, ES-C543-aplazar, ES-C543-adelantar, ES-C543-permanecer, ES-R543-repaso-536-543

\chapter{Más tarde de lo previsto}
\label{ch:es-mas-tarde-de-lo-previsto}
\hlchaptermodality{543}

\begin{chapteropening}
\textbf{By the end of this chapter:} \emph{I can say that something was resumed, delayed, put off, moved forward or stayed the same.}

Complete the last lesson of chapter 543: I can say that something was resumed, delayed, put off, moved forward or stayed the same.
\end{chapteropening}

\section[reanudar]{reanudar — to resume, to start again}

\label{lesson:ES-C543-reanudar}



\begin{quote}
Before the new one: say the Spanish for the start, then the Spanish for to end.
\end{quote}

\subsection*{You'll want to know: reanudar}
\textbf{reanudar} — \textquotedblleft{}to resume, to start again\textquotedblright{}.

\begin{quote}
Después de la lluvia reanudaron el partido.
\end{quote}

\begin{quote}
\emph{After the rain they resumed the match.}
\end{quote}

\emph{Reanudarse} — to start up again: \emph{Se reanudó el servicio.}

\begin{cousinweb}
\emph{reanudar} is built on \emph{anudar} 'to tie a knot' with \emph{re-} 'again'. \emph{Anudar}, older \emph{añudar} (EtymDB), is built on \emph{nudo} 'a knot', from Latin \emph{nōdus} 'a knot' (Lewis \& Short). To resume is to tie the thread again.
\end{cousinweb}

\begin{grammarlens}[title={What you've built}]
One more word for telling, explaining and describing at B1.
\end{grammarlens}

\subsection*{Your turn}
\begin{itemize}
\raggedright
  \item \emph{Say it:} \emph{reanudar}
  \item \emph{Say it:} \emph{reanudar}, once more
  \item \emph{Say it:} a sentence of your own with \emph{reanudar}
  \item \emph{Recall:} say the Spanish for the start, then the Spanish for to end, then say \emph{reanudar} again
\end{itemize}

\begin{tcolorbox}[breakable,colback=teal!4,colframe=teal!35!black,title={Before you move on}]
What does reanudar mean? (\textquotedblleft{}To resume, to start again\textquotedblright{}.) Say it once more.
\end{tcolorbox}
\section[retrasar]{retrasar — to delay, to put back}

\label{lesson:ES-C543-retrasar}



\begin{quote}
Before the new one: say the Spanish for to end, then the Spanish for to resume.
\end{quote}

\subsection*{You'll want to know: retrasar}
\textbf{retrasar} — \textquotedblleft{}to delay, to put back\textquotedblright{}.

\begin{quote}
La nieve retrasó el vuelo.
\end{quote}

\begin{quote}
\emph{The snow delayed the flight.}
\end{quote}

\emph{Retrasarse} is to be late: \emph{El tren se retrasó.}

\begin{cousinweb}
\emph{retrasar} is usually explained as built on \emph{tras} 'behind', which comes from Latin \emph{trāns} 'across, beyond' (Lewis \& Short), with \emph{re-} and the verb ending \emph{-ar}.
\end{cousinweb}

\begin{grammarlens}[title={What you've built}]
One more word for telling, explaining and describing at B1.
\end{grammarlens}

\subsection*{Your turn}
\begin{itemize}
\raggedright
  \item \emph{Say it:} \emph{retrasar}
  \item \emph{Say it:} \emph{retrasar}, once more
  \item \emph{Say it:} a sentence of your own with \emph{retrasar}
  \item \emph{Recall:} say the Spanish for to end, then the Spanish for to resume, then say \emph{retrasar} again
\end{itemize}

\begin{tcolorbox}[breakable,colback=teal!4,colframe=teal!35!black,title={Before you move on}]
What does retrasar mean? (\textquotedblleft{}To delay, to put back\textquotedblright{}.) Say it once more.
\end{tcolorbox}
\section[aplazar]{aplazar — to postpone, to put off}

\label{lesson:ES-C543-aplazar}



\begin{quote}
Before the new one: say the Spanish for to resume, then the Spanish for to delay.
\end{quote}

\subsection*{You'll want to know: aplazar}
\textbf{aplazar} — \textquotedblleft{}to postpone, to put off\textquotedblright{}.

\begin{quote}
Aplazaron el partido por la lluvia.
\end{quote}

\begin{quote}
\emph{They postponed the match because of the rain.}
\end{quote}

It takes \emph{c} before \emph{e}: \emph{aplacé}. \emph{Aplazar algo hasta el lunes} — to put something off until Monday.

\begin{cousinweb}
\emph{aplazar} is usually explained as built on \emph{plazo} 'a time limit' with \emph{a-} and \emph{-ar}. \emph{Plazo} comes from Latin \emph{placitum} 'what has been agreed', from \emph{placēre} 'to please' (Lewis \& Short).
\end{cousinweb}

\begin{grammarlens}[title={What you've built}]
One more word for telling, explaining and describing at B1.
\end{grammarlens}

\subsection*{Your turn}
\begin{itemize}
\raggedright
  \item \emph{Say it:} \emph{aplazar}
  \item \emph{Say it:} \emph{aplazar}, once more
  \item \emph{Say it:} a sentence of your own with \emph{aplazar}
  \item \emph{Recall:} say the Spanish for to resume, then the Spanish for to delay, then say \emph{aplazar} again
\end{itemize}

\begin{tcolorbox}[breakable,colback=teal!4,colframe=teal!35!black,title={Before you move on}]
What does aplazar mean? (\textquotedblleft{}To postpone, to put off\textquotedblright{}.) Say it once more.
\end{tcolorbox}
\section[adelantar]{adelantar — to bring forward; to overtake}

\label{lesson:ES-C543-adelantar}



\begin{quote}
Before the new one: say the Spanish for to delay, then the Spanish for to postpone.
\end{quote}

\subsection*{You'll want to know: adelantar}
\textbf{adelantar} — \textquotedblleft{}to bring forward; to overtake\textquotedblright{}.

\begin{quote}
Adelantaron la reunión al lunes.
\end{quote}

\begin{quote}
\emph{They brought the meeting forward to Monday.}
\end{quote}

On the road it means to overtake. \emph{Adelantarse} is to get ahead.

\begin{cousinweb}
\emph{adelantar} is built on \emph{adelante} 'forward' (EtymDB), which comes from \emph{delante}, old Spanish \emph{denante}, built on Latin \emph{dē}, \emph{in} and \emph{ante} 'before'.
\end{cousinweb}

\begin{grammarlens}[title={What you've built}]
One more word for telling, explaining and describing at B1.
\end{grammarlens}

\subsection*{Your turn}
\begin{itemize}
\raggedright
  \item \emph{Say it:} \emph{adelantar}
  \item \emph{Say it:} \emph{adelantar}, once more
  \item \emph{Say it:} a sentence of your own with \emph{adelantar}
  \item \emph{Recall:} say the Spanish for to delay, then the Spanish for to postpone, then say \emph{adelantar} again
\end{itemize}

\begin{tcolorbox}[breakable,colback=teal!4,colframe=teal!35!black,title={Before you move on}]
What does adelantar mean? (\textquotedblleft{}To bring forward; to overtake\textquotedblright{}.) Say it once more.
\end{tcolorbox}
\section[permanecer]{permanecer — to remain, to stay}

\label{lesson:ES-C543-permanecer}



\begin{quote}
Before the new one: say the Spanish for to postpone, then the Spanish for to bring forward.
\end{quote}

\subsection*{You'll want to know: permanecer}
\textbf{permanecer} — \textquotedblleft{}to remain, to stay\textquotedblright{}.

\begin{quote}
Permaneció en el hospital una semana.
\end{quote}

\begin{quote}
\emph{He remained in hospital for a week.}
\end{quote}

Formal; in speech \emph{quedarse}. It takes a \emph{z}: \emph{permanezco}.

\begin{cousinweb}
\emph{permanecer} comes from Latin \emph{permanēre} 'to stay to the end, to last' (EtymDB), built on \emph{per} and \emph{manēre} 'to stay', with the Spanish ending \emph{-ecer}.
\end{cousinweb}

\begin{grammarlens}[title={What you've built}]
One more word for telling, explaining and describing at B1.
\end{grammarlens}

\subsection*{Your turn}
\begin{itemize}
\raggedright
  \item \emph{Say it:} \emph{permanecer}
  \item \emph{Say it:} \emph{permanecer}, once more
  \item \emph{Say it:} a sentence of your own with \emph{permanecer}
  \item \emph{Recall:} say the Spanish for to postpone, then the Spanish for to bring forward, then say \emph{permanecer} again
\end{itemize}

\begin{tcolorbox}[breakable,colback=teal!4,colframe=teal!35!black,title={Before you move on}]
What does permanecer mean? (\textquotedblleft{}To remain, to stay\textquotedblright{}.) Say it once more.
\end{tcolorbox}
\section[(dialogue)]{Repaso de los capítulos 536–543}

\label{lesson:ES-R543-repaso-536-543}



\begin{quote}
Say the Spanish for meanwhile and at the same time.
\end{quote}

\subsection*{Your turn}
Say these aloud:

\begin{itemize}
\raggedright
  \item meanwhile, at the same time, at once, as soon as, the moment that
  \item since then, until then, the next day, a short time ago, the day before yesterday
  \item previously, subsequently, beforehand, recently, lately
  \item the semester, the fortnight, the five-year period, the millennium, the period
  \item at the beginning of, in the middle of, at the end of, weekly, yearly
  \item the phase, the passing, the length, to go on, to go by
  \item to begin, to launch, the beginning, the start, to end
  \item to resume, to delay, to postpone, to bring forward, to remain
\end{itemize}

\begin{tcolorbox}[breakable,colback=teal!4,colframe=teal!35!black,title={Before you move on}]
Meanwhile? (\textbf{mientras tanto}.) At the same time? (\textbf{al mismo tiempo}.) At once? (\textbf{a la vez}.) As soon as? (\textbf{en cuanto}.) The moment that? (\textbf{tan pronto como}.) Since then? (\textbf{desde entonces}.) Until then? (\textbf{hasta entonces}.) The next day? (\textbf{al día siguiente}.) A short time ago? (\textbf{hace poco}.) The day before yesterday? (\textbf{anteayer}.) Previously? (\textbf{anteriormente}.) Subsequently? (\textbf{posteriormente}.) Beforehand? (\textbf{previamente}.) Recently? (\textbf{recientemente}.) Lately? (\textbf{últimamente}.) The semester? (\textbf{el semestre}.) The fortnight? (\textbf{la quincena}.) The five-year period? (\textbf{el lustro}.) The millennium? (\textbf{el milenio}.) The period? (\textbf{el período}.) At the beginning of? (\textbf{a principios de}.) In the middle of? (\textbf{a mediados de}.) At the end of? (\textbf{a finales de}.) Weekly? (\textbf{semanal}.) Yearly? (\textbf{anual}.) The phase? (\textbf{la fase}.) The passing? (\textbf{el transcurso}.) The length? (\textbf{la duración}.) To go on? (\textbf{prolongarse}.) To go by? (\textbf{transcurrir}.) To begin? (\textbf{comenzar}.) To launch? (\textbf{iniciar}.) The beginning? (\textbf{el comienzo}.) The start? (\textbf{el arranque}.) To end? (\textbf{finalizar}.) To resume? (\textbf{reanudar}.) To delay? (\textbf{retrasar}.) To postpone? (\textbf{aplazar}.) To bring forward? (\textbf{adelantar}.) To remain? (\textbf{permanecer}.)
\end{tcolorbox}
